% 原文 pp.12–20（印刷 pp.1130–1138）
% 第 1 讲第 2 部分：承接第 1 部分（止于公式 (1.8)）；含原文 §1.2 的续论及 §1.3–§1.6 的开头。

% ===== 原文 p.12 =====
其中 $U$ 是前面那个四次势。与以前一样，我们有两个平衡态（equilibrium state）$a$、$-a$。

现在从量子力学的角度考察这个理论。这时可以看出，仍然没有对称性破缺。看清这一点最简单的办法，是先考虑离散的空间 $M$。此时算符代数（operator algebra）是对应于 $M$ 各点的 Heisenberg 代数的有限张量积（tensor product），因此它本身就是一个 Heisenberg 代数；于是我们在量子力学情形用过的表示论（representation theory）论证表明，对称性破缺不会发生。

% [Tutor Note: 「有限张量积仍是 Heisenberg 代数」一句背后的定理是 Stone–von Neumann 唯一性：$M$ 是有限点集，故这里的算符代数只有有限个自由度，其指数化（Weyl 代数）在等价意义下只有唯一的不可约表示，于是任何对称性都可由该表示上的（投影）酉算符实现，Hilbert 空间不会劈成互不相通的部分。作者在 §1.3 指出该论证在无穷体积下失效，原因正在这里：无穷张量积 $\bigotimes_{x\in M}$ 不再有唯一的不可约表示（von Neumann 的无穷张量积理论），不同表示对应不同的「宏观真空」，对称性破缺才成为可能。]

不过，用路径积分（path integral）论证更有启发性。与前面一样，先假设对称性破缺发生。于是我们有两个实现（realization）$\mathcal{H}_+$、$\mathcal{H}_-$，以及两个正交的量子真空 $\Omega_\pm\in\mathcal{H}=\mathcal{H}_+\oplus\mathcal{H}_-$；当 $\hbar$ 很小时，它们在如下意义下「局域化」于平衡点 $a$、$-a$ 附近：
\[
\langle\Omega_\pm,\phi(x_1)\cdots\phi(x_n)\Omega_\pm\rangle\to(\pm a)^n,\qquad \hbar\to0 .
\]

与量子力学的情形一样，内积 $(\Omega_+,e^{-tH/\hbar}\Omega_-)$ 为零。另一方面，用 Feynman–Kac 公式计算它，我们将得到
\begin{equation}
  \tag{1.9}
  \label{eq:1-9}
  (\Omega_+,e^{-tH/\hbar}\Omega_-)\sim e^{-S_*(t)V/\hbar},
\end{equation}
其中 $S_*(t)$ 如前，$V$ 是空间的体积。理由是：最小作用量在空间无关的经典解 $\phi^M_*(\mu,\tau)=\phi_*(\tau)$（$\mu\in M$，$\tau\in\mathbb{R}$）上达到。

% [Tutor Note: (1.9) 的指数里出现体积因子 $V$，这是全段的关键：因为 $M$ 是紧的且作用量是空间积分，最省作用量的构型在空间上必须是常数，于是作用量恰为「单位体积的作用量 $S_*(t)$」乘 $V$。也正因为 $V$ 有限，$e^{-S_*(t)V/\hbar}$ 才非零，才能与「内积为零」矛盾；一旦 $V\to\infty$ 这个矛盾立刻消失——这正是下一节的分水岭。复算时注意 $S_*(t)$ 是 $M$ 为单位体积时的量（与 (1.7) 的定义一致），不要把它再乘一遍体积。]

由于 (1.9) 非零（这里体积 $V$ 有限是本质的），我们得到矛盾。

与量子力学中一样，我们可以定义两组关联函数（correlation function）$W^n_+$、$W^n_-$，它们由在最低能量点 $a$、$-a$ 附近的微扰论计算得到（当然，只有在可重整化的情形，即四维及以下，才能定义它们；此外还应记住，重整化并不是唯一确定的）。我们的论证表明，$W^n_+$、$W^n_-$ 并不是我们的量子理论某个实现的关联函数的小 $\hbar$ 展开。另一方面，函数
\[
W^n_0=\tfrac12\bigl(W^n_++W^n_-\bigr)
\]
则有机会成为实际（非微扰）关联函数的渐近展开。

当 $\dim(M)=1,2$ 时，量子理论 $\mathcal{H}$ 的存在性已在构造性场论（constructive field theory）中建立。在此情形可以证明：（在有限体积下）存在唯一的基态 $\Omega_0$，它是 $G$-不变的；并且该理论关于 $\Omega_0$ 的关联函数确实具有由 $W^n_0$ 给出的渐近展开。

% 原文 §1.3
\subsection{无限体积中量子场论的对称性破缺}

当空间 $M$ 的体积 $V$ 趋于无穷时，上一节的论证失效。表示论论证失效（即使存在紫外截断也失效），因为我们在有限体积情形所用的算符代数的正则表示，现在是对应于各点的空间的无穷张量积；于是我们必须赋予它意义，而可能并不存在 $G$-不变的做法。路径积分论证同样失效。确实，(1.9) 的右端为零，因而对 $(\Omega_+,e^{-tH/\hbar}\Omega_-)$ 的两种计算给出相同答案，我们推不出任何矛盾。此外，若假定我们有一个算符代数的表示 $\mathcal{H}$，其中含两个真空 $\Omega_+$、$\Omega_-$，那么

% ===== 原文 p.13 =====
利用公式
\begin{equation}
  \tag{1.10}
  \label{eq:1-10}
  \langle\Omega_+,\phi(x_1^*)\cdots\phi(x_n^*)\Omega_-\rangle
  =\int_{\substack{\phi:X\to\mathbb{R},\ \phi\to\pm a,\ t\to\pm\infty}}
  \phi(x_1)\cdots\phi(x_n)\,e^{-S(\phi)/\hbar}D\phi,
\end{equation}
（其中对 $x=(\mu,t)$ 有 $x^*=(\mu,it)$），我们推断：对任意的 $x_1,\dots,x_n$，内积 (1.10) 为零。这表明空间 $\mathcal{H}$ 分裂为正交直和
\[
\mathcal{H}=\mathcal{H}_+\oplus\mathcal{H}_-,
\]
其中 $\mathcal{H}_\pm$ 是由真空 $\Omega_+$、$\Omega_-$ 生成的、彼此独立的两个实现的 Hilbert 空间。这意味着对称性破缺可以发生，正如经典理论中那样：量子效应并不必然强到足以恢复对称性。

% [Tutor Note: (1.10) 为何恒为零，作者没有明说，其依据是无穷体积：边界条件 $\phi\to\pm a$（$t\to\pm\infty$）要求场在整个空间上从 $-a$ 过渡到 $+a$，最省作用量的这种构型是一个充满全空间的「畴壁」，其作用量正比于 $V\to\infty$，故右端按 $e^{-\mathrm{const}\cdot V}$ 熄灭。注意这与 (1.9) 恰好相反：那里 $V$ 有限故指数非零，这里 $V=\infty$ 故指数为零。另需留意 $x^*=(\mu,it)$ 表示插入点取在欧氏点上（Wick 转动后的时间），这是把 $\Omega_+$、$\Omega_-$ 之间的矩阵元写成路径积分的必要步骤；$D\phi$ 未写归一化因子，它被 $1=\langle\Omega_\pm,\Omega_\pm\rangle$ 的约定吸收。]

这正是量子场论中实际发生的事。更精确地说，情况如下。

\begin{enumerate}
  \item[(i)] 若对称群 $G$ 是有限群，则从时空维数 2 起，对称性就可以在无穷体积中破缺。例如，在我们考察的那个例子里，对称性破缺确实发生。
  \item[(ii)] 若对称群 $G$ 是连通 Lie 群，则从时空维数 3 起，对称性破缺才可能发生。它在维数 2 不发生，这是一个非凡而极其重要的事实，我们将在 §1.5 讨论。该事实由 S. Coleman 于 1973 年证明（文献见本讲末尾）。
\end{enumerate}

因此，对称性破缺是一种「红外」效应，它与理论在大距离上的行为相联系，并且对时空维数极为敏感。

\textbf{注（Remark）.}
在某种意义下，在维数 2，一旦我们考虑在不同真空之间插值的孤子（soliton），所有真空就都成为一个更统一的故事的一部分。这一点在第 5 讲末尾的一个例子中作了说明；孤子在第九讲中有更系统的讨论。

% 原文 §1.4
\subsection{关联函数的无穷体积渐近}

让我们看看，在有限体积量子场论的无穷体积极限中，对称性是如何破缺的。

考虑从有限体积出发、逼近无穷体积极限的任意一种方式。例如，我们可以把时空取为紧 $d$-流形 $M$，并把 $M$ 的度量 $g$ 放大为 $g\to tg$，其中 $t\to\infty$。在此过程中，$M$ 的任一局部区域都收敛到 $\mathbb{R}^d$。我们考虑在点 $x_i$ 处插入的 $\mathcal{O}_i(x_i)$（其中 $\mathcal{O}_i$ 是任意一组局域算符，不必是基本场），这些点在 $t\to\infty$ 时收敛到 $\mathbb{R}^d$ 中指定的点。\footnote{在我们将于第 19 讲做的、对 Donaldson 理论的一个应用中，除了局域于点的算符之外，还必须考虑与曲面 $\Sigma_\alpha$ 相联系的算符。但那些算符不过是通常局域算符在 $\Sigma_\alpha$ 上的积分；在应用我们即将给出的讨论时，可以把 $\mathcal{O}_i$ 取为这些后者的局域算符。} 然后我们设法分析关联函数
\begin{equation}
  \tag{1.11}
  \label{eq:1-11}
  \langle\mathcal{O}_1(x_1)\cdots\mathcal{O}_n(x_n)\rangle_t
\end{equation}
——即给定乘积在给定 $t$ 值下的关联函数——在大 $t$ 极限下的行为。

另一个本质上是同一件事的例子是：可以把 $M$ 取为 $\mathbb{R}^d$ 中半径为 $t$ 的球，并在球的边界上取某种边界条件。

% ===== 原文 p.14 =====
（更一般地，$M$ 可以是 $\mathbb{R}^d$ 中任意区域，其尺寸被放大 $t$ 倍。）然后人们考虑这样的关联函数：它们通过在服从边界条件的场 $\phi$ 上做路径积分来定义。在此情形下，我们将用符号 (1.11) 来描述在点 $x_i\in\mathbb{R}^d$ 处固定的那些算符的关联函数，这些点在 $t\to\infty$ 时保持不变。

我们首先断言：若乘上一个适当的函数 $C(t)$，则 (1.11) 中的关联函数在 $t\to\infty$ 时具有极限。函数 $C(t)$ 依赖于第一种情形中带度量的 $M$ 的选择细节，或第二种情形中区域的形状与边界条件。一旦关联函数乘上 $C(t)$ 而具有极限，关于该极限我们有如下断言。设 $\langle\mathcal{O}_1(x_1)\cdots\mathcal{O}_n(x_n)\rangle_{(i)}$ 是理论在第 $i$ 个真空中的相应关联函数。则断言是
\begin{equation}
  \tag{1.12}
  \label{eq:1-12}
  \lim_{t\to\infty}C(t)\langle\mathcal{O}_1(x_1)\cdots\mathcal{O}_n(x_n)\rangle_t
  =\sum_i p_i\,\langle\mathcal{O}_1(x_1)\cdots\mathcal{O}_n(x_n)\rangle_{(i)}
\end{equation}
其中 $p_i$ 为某些复数。

直观上，这个重要公式的理由是：当 $t\to\infty$ 时，为了避免作用量过大，空间几乎处处都极其接近某一个真空。但这个充满空间的真空可以是任意一个，即对任意 $i$ 的第 $i$ 个真空。于是关联函数被写成第 $i$ 个真空中关联函数对 $i$ 的求和。不同真空出现不同的常数 $p_i$，是因为函数 $C(t)$ 当 $t\to\infty$ 时的行为，自然只确定到乘一个复常数。没有理由认为对每个 $i$ 都会自然地出现「同一个」常数，所以我们必须用不同的常数去加权不同的真空。（某些 $p_i$ 可以为零，如果在 $t\to\infty$ 极限下某些真空的贡献相比另一些可以忽略。）

在精神上——尽管也许难以严格化——对 (1.12) 的一个有吸引力的解释如下。本理论的关联函数 $\langle\mathcal{O}_1(x_1)\cdots\mathcal{O}_n(x_n)\rangle$ 服从大量由算符乘积展开（OPE）得出的线性关系。由于这些关系是线性的，解空间是一个向量空间。我们有理由希望：这个向量空间的一组基由各真空中的关联函数 $\langle\mathcal{O}_1(x_1)\cdots\mathcal{O}_n(x_n)\rangle_{(i)}$ 给出。(1.12) 的左端定义了此向量空间中的一个元素（因为左端的关联函数当然服从由算符乘积展开得出的一切关系）。所以它必定具有 (1.12) 中所写的那种展开。

% [Tutor Note: 这一段是全节唯一未证明的断言：「OPE 线性关系的解空间由各真空中的关联函数张成」只是一句 "one can reasonably hope"。它的严格化需要两步：其一，用 Osterwalder–Schrader 重构把关联函数翻译成 Hilbert 空间上的真空态与算符代数；其二，证明在 $t\to\infty$ 极限下态空间按「宏观真空」分解（即超选择分块的分解，等价于对 $C^*$ 代数做中心分解）。注意 $p_i$ 是复数而不必是正实数，所以 (1.12) 不是概率式的凸组合；在反射正定的欧氏理论中若取 $C(t)>0$，通常可期望 $p_i\ge0$，但作者有意不作此假设。]

注意，若不止一个 $p_i$ 非零，则由从紧流形 $M$（带或不带边界）出发并取 $t\to\infty$ 得到的关联函数不服从集团分解（cluster decomposition）。除此之外，它们服从场论的所有公理。从公理化的观点看（回忆 [I-Axiomatics] 第 2 讲），人们可以希望证明：服从关联函数的一切公理、只不服从集团分解的函数，可以按 (1.12) 用确实服从集团分解的符号展开。

% [Tutor Note: 指本丛书第 I 卷中的综述 [I-Axiomatics]（公理化场论），其第 2 讲讨论 Wightman 公理与 Osterwalder–Schrader 公理；「除集团分解外满足一切公理」的说法即指：这些函数仍是某个 Hilbert 空间上的真空期望值，只是该空间不是真空唯一的。]

现在我们举出一些可以确定 $p_i$ 的简单情形。若 $M$ 是带平坦度量的环面，则在大多数简单情形下 $p_i$ 相等，因为 $M$ 上的路径积分具有迹的解释，而给定能量的所有态对迹的贡献相同。（在某些条件下可以避免这种相等，例如存在未破缺的离散规范对称性时；但这是我今天不打算讲的内容。）作为另一类例子，设 $M$ 是

% ===== 原文 p.15 =====
$\mathbb{R}^d$ 中的球，并在球面上给定某些边界条件。设我们有一个弱耦合的标量场 $\phi$ 的理论，它具有对称性 $\phi\to-\phi$，以及两个经典基态 $\phi=a$ 与 $-a$。把相应的两个量子真空标为 $+$ 与 $-$。若取边界条件为边界上 $\phi=a$，则当 $t\to\infty$ 时，除了 $+$ 真空之外越来越难看到别的东西，于是 $p_+\neq0$，$p_-=0$。（适当选取 $C(t)$ 可使 $p_+=1$。）反之，若边界上 $\phi=-a$，则 $p_+=0$，$p_-\neq0$。最后，可以取边界上 $\phi=0$。这保持了交换两个真空的那种对称性，故 $p_+=p_-$。

% 原文 §1.5
\subsection{连续对称性破缺}

现在我们来考虑量子场论中的对称性破缺，其中对称群是一个连通 Lie 群。典型的例子是复值玻色场 $\phi$，以及作用量
\begin{equation}
  \tag{1.13}
  \label{eq:1-13}
  \mathcal{L}(\phi)=\int d^dx\left(\frac12|\nabla\phi|^2+\bigl(|\phi|^2-a^2\bigr)^2\right).
\end{equation}

% [Tutor Note: 这里 $\mathcal{L}(\phi)$ 带有 $\int d^dx$，是作用量（action）而非拉氏量密度；本书全篇用 $\mathcal{L}$ 表示作用量，与 §1.6 中「$\mathcal{L}=\mathcal{L}(\phi)$，$X$ 为 Euler–Lagrange 方程的解空间」的用法一致，读文献时不要与通常的 $\mathcal{L}$（密度）混淆。其次，$\nabla$ 必须是 $d$ 维（欧氏）梯度而非仅空间梯度，否则作用量中已含时间积分会使理论失去 Lorentz 不变性。第三，势 $(|\phi|^2-a^2)^2$ 前没有独立的耦合常数，系数取 1 等价于用 $\phi$ 的归一化把耦合固定下来；若要真正做弱耦合展开，须另设小参数（例如把势写成 $\lambda(|\phi|^2-a^2)^2$ 并取 $\lambda\to0$）。]

此作用量具有 $U(1)$ 对称性，作用为 $\phi\to e^{i\theta}\phi$。在经典理论中，最低能量态是 $\phi=ae^{i\theta}$，$\theta\in[0,2\pi)$，所以有对称性破缺。我们将设法弄清量子理论中是否也存在对称性破缺。

上面我们证明了，在一维情形（量子力学）不发生对称性破缺。我们是对实玻色子 $\phi$ 做的；但对复玻色子，论证（用路径积分）甚至更好。例如，对作用量 (1.13)，能量的经典极小点集合是圆 $|\phi|=a$，它是连通的。因此，从一个经典极小点 $ae^{i\theta_1}$ 到另一个 $ae^{i\theta_2}$，不需要翻越势的「驼峰」，而可以沿极小圆走，于是所需作用量比先前更小。因此，上面描述的路径积分计算会表明：在复情形下，态 $\psi_+$、$\psi_-$ 之间的联系比实情形更强。

% [Tutor Note: 这里的 $\psi_\pm$ 是上一部分量子力学双阱讨论中「局域化于 $\pm a$ 附近」的两个态（作者在 (1.7) 之后特别说明：论证只用到它们局域化于经典平衡态，并未用到它们是 $\mathbb{R}$ 上的函数），并非场论中的真空 $\Omega_\pm$。二者的关键差别在指数上：$\psi_\pm$ 之间的隧穿振幅 $\sim e^{-S_0/\hbar}$ 与空间体积无关，故量子力学里隧穿总能恢复对称性；而场论中 $\Omega_\pm$ 之间的矩阵元按 (1.9) 为 $e^{-S_*(t)V/\hbar}$，随 $V\to\infty$ 熄灭，所以场论中隧穿远不足以恢复对称性。连续情形下真正起决定作用的也不是隧穿，而是红外涨落（Coleman 定理；Mermin–Wagner 型论证），这正是下文转向 $\sigma$ 模型与 Goldstone 定理的原因。]

这一效应也许会使我们认为，场论中连续对称性破缺或许不像离散对称性破缺那样容易发生。而确实，结果表明连续对称性破缺在二维不能发生，只能在维数 $\ge3$ 发生。

遗憾的是，这样简单的路径积分论证太粗糙，不足以证明连续对称性破缺在二维不发生。确实，容易看出：在计算 $(\Omega_+,e^{-tH}\Omega_-)$ 的积分中，所有路径在无穷体积下都具有无穷作用量，所以我们推不出任何矛盾。因此，让我们通过考察最简单的例子，来说明为什么对称性在二维得以保持、而在二维以上破缺。

最简单的例子是 $d$ 维自由无质量实标量玻色场理论，它由作用量
\begin{equation}
  \tag{1.14}
  \label{eq:1-14}
  \mathcal{L}_0(\phi)=\int d^dx\left(\frac12(\nabla\phi)^2\right)
\end{equation}
定义。

这个理论具有平移对称性 $\phi\to\phi+c$。让我们证明：该对称性在 $d>2$ 时破缺，在 $d=2$ 时保持。

% ===== 原文 p.16 =====
对 $d>2$，对称性破缺是显然的。确实，在此情形我们有一个由基本场 $\phi$ 生成的 Wightman 场论，它满足 Wightman 公理（[I-Axiomatics] 第 1 讲）。在此理论中，算符 $\phi$ 的单点函数为零，而算符 $\phi+c$ 的单点函数为 $c$。因此变换 $\phi\to\phi+c$ 不保持单点函数。这表明在 $d>2$ 情形下，我们实际上并没有单一的实现，而是有一族实现 $\mathcal{H}_c$，由直线上的点 $c$ 参数化。理论 $\mathcal{H}_c$ 由条件 $\langle\Omega_c,\phi\Omega_c\rangle=c$ 定义，其中 $\Omega_c$ 是 $\mathcal{H}_c$ 的真空。真空 $\Omega_c$「局域化」于经典平衡态 $c$ 附近。映射 $\phi\to\phi+c$ 把 $\mathcal{H}_b$ 变为 $\mathcal{H}_{b+c}$。

% [Tutor Note: 「显然」一词省掉了两点。第一，$\langle\Omega,\phi\Omega\rangle=0$ 的根据是 (1.14) 具有 $\phi\to-\phi$ 的 $\mathbb{Z}_2$ 对称性（作用量只含 $(\nabla\phi)^2$），而 $\phi+c$ 破坏了它；仅凭平移不变性只能断言该单点函数是常数，不能断言它为零。第二，严格说 $\phi+c$ 属于「另一个」理论：Wightman 重构定理保证 $\phi+c$ 仍生成一套满足公理的关联函数，但其真空期望值为 $c\neq0$，故给出与原来不等价的表示 $\mathcal{H}_c$。这两步都依赖 $\phi$ 本身作为算符值分布存在，而这恰是 $d>2$ 才成立的（见下一段 $d=2$ 的反例），所以论证不能平移到二维。]

对 $d=2$，情形不同。问题在于：算符 $\phi$ 在二维没有定义，尽管它的导数有定义。确实，在二维我们有
\begin{equation}
  \tag{1.15}
  \label{eq:1-15}
  \langle\partial\phi(x)\partial\phi(y)\rangle=-\partial_x\partial_y\ln|x-y|,
\end{equation}
于是，如果算符 $\phi$ 以某种方式有定义，我们就会有
\begin{equation}
  \tag{1.16}
  \label{eq:1-16}
  \langle\phi(x)\phi(y)\rangle=-\ln|x-y|+C,
\end{equation}
这与正定性（positivity）矛盾（(1.16) 右端的函数不是正定的）。

% [Tutor Note: (1.15)–(1.16) 与 (1.14) 的归一化相差一个常数：按 $\mathcal{L}_0=\frac12\int d^2x(\nabla\phi)^2$，二维传播子由 $-\Delta G=\delta^{(2)}$ 定出为 $G(x-y)=-\frac{1}{2\pi}\ln|x-y|$，即 (1.16) 右端应带系数 $-\frac{1}{2\pi}$。作者显然取了使系数为 1 的归一化（等价于对 $\phi$ 重新标度并使 $\lambda$ 随之改标）。这不影响「违反正定性」的结论，但下一节 $|x|^{-k^2/\lambda^2}$ 中的指数会因此不带 $2\pi$；读者用标准归一化复算时，应得到 $|x|^{-k^2/(2\pi\lambda^2)}$。至于「违反正定性」：$-\ln|x-y|$ 在 $|x-y|\to\infty$ 时无下界，而正定性要求 $\sum_{i,j}c_ic_j\langle\phi(x_i)\phi(x_j)\rangle\ge0$，取相距很远的两点即可看出矛盾。]

让我们把这一点说得更精确。在二维我们所拥有的，是由从 Schwartz 函数到算符的 Wightman 映射 $\phi$ 生成的量子场论；它并非定义在整个 Schwartz 函数空间 $S(V)$ 上，而只定义在 $V$ 上积分为零的 Schwartz 函数空间 $S_0(V)$ 上。于是 $\phi$ 的导数可以通过
\[
\partial_i\phi(f)=-\phi(\partial_i f)
\]
定义在所有 Schwartz 函数上（右端有意义，因为 $\int\partial_i f=0$）。

然而，在由这样的 $\phi$ 定义的理论中，对称性破缺的问题根本不会出现，因为一开始就没有对称性：在空间 $S_0(V)$ 上，映射 $\phi$ 与 $\phi+c$ 是同一个。所以，为了提出关于对称性破缺的问题，我们应当考虑我们的算符代数的一个扩张，它将具有非平凡的对称性作用。这样做最合理的办法如下。

我们不考察 $\mathbb{R}$ 值无质量标量 $\phi$ 的理论，而考察取值于圆、且作用量相同的场 $\phi$ 的理论。取圆为 $S_\lambda=\mathbb{R}/2\pi\lambda\mathbb{Z}$。作用量与前面相同。在此情形，局域泛函 $\phi$ 没有定义，但我们有局域泛函 $e^{ik\phi/\lambda}$，$k\in\mathbb{Z}$。此理论中的局域泛函是 $e^{i\phi/\lambda}$ 的 Laurent 多项式，其系数是场 $\phi$ 各导数的微分多项式，而不是 $\phi$ 本身的微分多项式。此理论中的平移对称性是 $U(1)$ 对称性：对满足 $|z|=1$ 的复数 $z$，有 $z\circ e^{ik\phi/\lambda}=z^k e^{ik\phi/\lambda}$，而 $z$ 平凡地作用于 $\phi$ 的导数。经典地，此理论有一族真空（平衡态），由 $\phi=c$，$c\in S_\lambda$ 给出。

我们要证明：在这个系统中，对称性在量子力学意义上不破缺。只需证明：对任意形如 $\mathcal{O}=P(\phi)e^{ik\phi/\lambda}$（$k\neq0$）的局域算符 $\mathcal{O}$，单点函数 $\langle\Omega,\mathcal{O}(0)\Omega\rangle$ 为零。\footnote{所有局域算符都是这种形式；并且可以用算符乘积展开证明：若存在一个并非不变的 $n$ 点函数，则存在这样的一个单点函数。} 让我们在

% ===== 原文 p.17 =====
$P=1$ 的情形证明（一般情形的证明类似）。证明：由自由理论中的算符乘积展开（参见 [I-Perturbative QFT] 第 3 讲）我们得到
\[
\langle\Omega,\mathcal{O}(x)\mathcal{O}^*(0)\Omega\rangle=|x|^{-k^2/\lambda^2},
\]
所以这个两点函数在无穷远处趋于零。但由集团性质，该函数在无穷远处的极限是 $|\langle\Omega,\mathcal{O}(0)\Omega\rangle|^2$。故 $\langle\Omega,\mathcal{O}(0)\Omega\rangle=0$。

看对称性不破缺的另一种方式如下。量子化理论的 Hilbert 空间形如
\[
\mathcal{H}=F\otimes F^*\otimes l_2(\mathbb{Z})=\bigoplus_{k\in\mathbb{Z}}(F\otimes F^*)_k,
\]
其中 $F$ 是 Fock 空间。对应于 $\phi$ 的导数的算符保持这个分解，而算符 $e^{ik\phi/\lambda}$ 把空间 $(F\otimes F^*)_{k'}$ 映到空间 $(F\otimes F^*)_{k+k'}$。真空向量 $\Omega$ 属于零分量 $(F\otimes F^*)_0$。这意味着此理论中算符的一切关联函数在 $U(1)$ 作用下不变。特别地，
\[
\langle e^{ik\phi/\lambda}\rangle=\delta_{0k},
\]
这表明真空 $\Omega$ 并不「局域化」于任何经典基态附近，而是均匀地「铺展」在经典基态空间 $S_\lambda$ 上。因此，$U(1)$ 对称性没有破缺。

% [Tutor Note: 这个 Fock 空间分解需要解释：圆值场 $\phi$ 的共轭动量按 $p=k/\lambda$（$k\in\mathbb{Z}$）量子化，故单粒子空间是 $\mathbb{Z}$-分次的；$\mathcal{H}=F\otimes F^*\otimes l_2(\mathbb{Z})$ 中的 $l_2(\mathbb{Z})$ 记录这个动量/绕数扇区，$F\otimes F^*$ 记录左右 mover 的激发。关键在于真空 $\Omega$ 落在 $k=0$ 分量，而 $\phi$ 本身不是良定义的算符（$e^{i\phi/\lambda}$ 才是，它把相邻 $k$ 分量混合），这正是作者前面强调「只有 $e^{ik\phi/\lambda}$ 是局域泛函」的算符代数表述。也正因如此，$\phi\to\phi+c$ 只有在 $c=2\pi\lambda n$ 时才是恒等映射；圆值理论里真正的对称性只剩靶空间的 $U(1)$ 转动，它在 $k\neq0$ 分量上给出相因子 $z^k$，其是否破缺恰由 $\langle e^{ik\phi/\lambda}\rangle=\delta_{0k}$ 判定。]

\textbf{注 1（Remark 1）.} 看一看为什么我们关于 $d>2$ 存在对称性破缺的论证在有限体积下失效，是有启发的；因为我们知道，有限体积下不应有对称性破缺。问题在于：在有限体积下，场 $\phi$ 没有定义，尽管它的导数有定义。确实，由于时空是时间直线与紧空间的乘积，在大距离上它看起来就只是一条直线；$\phi$ 的两点函数（即时空的 Green 函数）在 $|t|$ 大时看起来像一维 Green 函数，即 $-|t|+C+o(1)$。此函数没有下界，因而违反正定性。此外，由考察 $d=2$ 的情形可知，若除一个空间方向外其余方向都被紧化，情形也一样。

\textbf{注 2（Remark 2）.} 上面我们考察了维数 $d>2$ 的自由无质量标量 $\phi$ 的理论，发现其量子真空空间是 $\phi$ 的取值空间（即靶空间 $\mathbb{R}$）。更一般地，若考虑时空为 $\mathbb{R}^d$（$d\ge2$）、靶空间为 $M$（一个黎曼流形）的 $\sigma$ 模型（sigma model），则量子真空空间以及经典基态空间都将是 $M$。（这里你应当暂时忘记，这个非线性 $M$ 的 $\sigma$ 模型是不可重整的。在今天讲座的第 8 部分，我们将把这个 $\sigma$ 模型理解为红外有效理论。）特别地，作为黎曼流形的 $M$ 可以从量子理论中被还原出来，即作为量子真空的模空间（moduli space of vacua）（见下面的注 3）。

另一方面，我们看到对 $d=2$，取值于圆的场 $\phi$ 的理论（它等同于靶空间为 $S_\lambda$ 的二维 $\sigma$ 模型）只有一个真空 $\Omega$。一般二维 $\sigma$ 模型都是这样：其量子真空的模空间通常非常小，与经典基态空间（$=$ 靶空间）不重合。特别地，靶空间不能从量子理论内在地还原出来。例如，圆 $S_\lambda$ 不能从上面考虑的、映入 $S_\lambda$ 的映射的理论中被内在地还原出来。这其实有很好的理由——可以证明，附着于圆 $S_\lambda$ 与 $S_{1/\lambda}$ 的理论是等价的（例如它们的配分函数相同——见 [II-CFT] 第 1 讲）。这正是镜像对称（mirror symmetry）理论的出发点。

% ===== 原文 p.18 =====
\textbf{注 3（Remark 3）.} 考虑一个作用量为
\[
\mathcal{L}=\int d^dx\left(\frac12(\nabla\phi)^2+U(\phi)\right)
\]
的场论，其中 $\phi$ 取值于某个黎曼流形 $M$，$U$ 是 $M$ 上的势函数（$U\ge0$）。设 $M(0)$ 是 $U$ 的零点集。假定 $M(0)$ 非空且光滑，并假定 $U$ 的二阶导数 Hessian 矩阵在 $T_xM/T_xM(0)$ 上、对 $x\in M(0)$ 非退化。设有一个 Lie 群 $G$，它通过等距作用在 $M$ 上、保持 $U$ 不变，且在 $M_0$ 上可迁。在此情形可以证明（在物理的严格性水平上）：由 $\mathcal{L}$ 描述的理论的「红外行为」，与靶空间取为经典真空空间 $M(0)$ 的 $\sigma$ 模型的「红外行为」相同。这一陈述的精确含义在第 1.7 节中解释。（我们可以用截断路径积分来定义这些理论，从而回避不可重整性问题：在截断路径积分中，积分取在步长为 $\Lambda^{-1}$ 的格点上定义的场之上，或取在相对于某个坐标系只有 $|k|<\Lambda$ 的 Fourier 模式的场之上；在这种框架中，截断 $\Lambda$ 不趋于无穷。）这一事实可以启发式地解释：若一个函数 $\phi$ 要具有低能量，它必须缓慢振荡，并且还要靠近极小点轨迹 $M(0)$，即必须接近一个映入 $M(0)$ 的映射。所以低能物理与缓变的、映入 $M(0)$ 的映射有关。

% [Tutor Note: 此处原文作 "acts transitively on $M_0$"，而同段其余各处一律写 $M(0)$（$U$ 的零点集），应视为同一对象的笔误，译文按 $M(0)$ 处理。至于定理本身，作者说 "one can show (at the physical level of rigor)"，其证明思路是：把 $\phi$ 分解为沿 $M(0)$ 的零模（Goldstone 模）与横向的 massive 模，后者质量由 Hessian 的非退化性给出正的下界（这正是「Hessian 在 $T_xM/T_xM(0)$ 上非退化」这一条件的用处），在低能下把 massive 模积分掉只留下 $M(0)$-值映射的 $\sigma$ 模型；$G$ 在 $M(0)$ 上可迁则保证横向质量矩阵在 $M(0)$ 上处处同型。严格的讨论见 [I-Classical Fields] 与后续第 1.7 节；另注意 $U\ge0$ 与 $M(0)\neq\varnothing$ 保证了基态能量为零、真空模空间就是 $M(0)$。]

因此，连续对称性破缺作为红外效应，在由 $\mathcal{L}$ 描述的理论中发生，当且仅当它在相应的 $\sigma$ 模型中发生。于是，由注 2 可知，对称性破缺倾向于不在维数 2 发生，而倾向于在维数 $>2$ 发生。

然而，这些大多只是启发式论证。在下一节，我们将借助 Goldstone 定理，更系统地处理连续对称性破缺问题。

% 原文 §1.6
\subsection{Goldstone 定理}

回忆经典场论中 Noether 定理与守恒流的标准形式。设我们在 $d$ 维时空 $V$ 中有一个作用量 $\mathcal{L}=\mathcal{L}(\phi)$。把相应的 Euler–Lagrange 方程的解空间记为 $X$。

设 $G^s$ 是此作用量的单参数对称群。令
\[
D_\phi:=\left.\frac{d}{ds}\right|_{s=0}G^s\phi .
\]
我们假定 $D_\phi$ 是 $\phi$ 的局域泛函。

设 $\eta\in\Omega^1(X,\Omega^{d-1}(V))$ 是解空间上的典范一形式（canonical one-form），它在 [I-Classical Fields] 与 [I-Homework] 中讨论过（$X$ 上的典范二形式定义为 $\int_C d_X\eta$，其中 $C$ 是 $d-1$ 维闭链）。例如，对无质量标量的自由理论，$\eta$ 的公式为
\begin{equation}
  \tag{1.17}
  \label{eq:1-17}
  \eta(\delta\phi)(x)=\delta\phi(x)*d\phi(x).
\end{equation}

% [Tutor Note: [I-Classical Fields] 与 [I-Homework] 指本丛书第 I 卷中关于经典场论与习题的两部分；$\eta$ 的构造是：在解空间 $X$ 上，场变分 $\delta\phi$ 与 $\phi$ 的时空导数配对，$\eta(\delta\phi)=\delta\phi\,{*}d\phi$ 是把 $\delta\phi$ 与 $*d\phi$（一个 $(d-1)$-形式）相乘得到的 $(d-1)$-形式。这里 $*$ 是时空 $V$ 上的 Hodge 星号，注意它作用在 $d\phi$ 上而非 $\delta\phi$ 上；若把 $*$ 的位置搞错，$J=\eta(D_\phi)$ 会变成一个 $1$-形式而非 $(d-1)$-形式，后面的电荷积分 $\int_C J$（在 $d-1$ 维闭链上）就无法定义。]

设 $J\in\Omega^0(X,\Omega^{d-1}(V))$ 由公式 $J=\eta(D_\phi)$ 定义。由于 $D_\phi$ 是局域的，$J$ 也是局域的。于是 $J$ 是 $X$ 上取值于 $\Omega^{d-1}(V)$ 的局域泛函。例如，若 $\mathcal{L}$ 是无质量标量自由理论的作用量，且 $G^s\phi=\phi+s$，则 $D_\phi=1$，故 $J(x)=*d\phi(x)$。局域泛函 $J$ 称为对应于对称性 $G^s$ 的流（current）。流的主要性质是它守恒，即 $d_V J=0$。这由 $d_V\eta=0$ 这一事实推出。

% [Tutor Note: 原文此句印作 "This follows form the fact that $d_V\eta=0$"，"form" 系 "from" 之误（p.1136 末行）。另注意此处「守恒」是外微分意义下的 $d_VJ=0$，等价于 $\partial_\mu J^\mu=0$；而 $J$ 是 $(d-1)$-形式而非矢量，二者通过 Hodge 星号对应，做分量计算时别忘了 $J^\mu=(\ast J)$ 的指标位置。]

% ===== 原文 p.19 =====
在有限体积下，定义电荷泛函（charge functional）$Q=\int_C J(x)$ 是有用的，其中 $C$ 是某个类空闭链（例如 $t=\text{const}$）。由于流守恒，只要 $C$ 代表「空间」的基本同调类，这个量就与闭链的选取无关。若 $\Omega=\int_C d_X\eta$ 是 $X$ 上的非退化二形式，则 $Q$ 是定义对称群 $G^s$ 的 Hamiltonian，其意义是
\[
\left.\frac{d}{ds}\right|_{s=0}(G^s)^*F=\{F,Q\},
\]
其中 $\{\,,\}$ 是 $S$ 上的 Poisson 括号，$F$ 是 $X$ 上任一局域泛函。

在无穷体积下，泛函 $Q$ 不一定有定义，因为积分不收敛。此时方便的做法是选取时间坐标 $t$，令 $C=C_0=\{t=0\}$ 为初值面，然后定义「截断电荷泛函」
\begin{equation}
  \tag{1.18}
  \label{eq:1-18}
  Q_f=\int_{C_0}f(x)J(x),
\end{equation}
其中 $f:C_0\to\mathbb{R}$ 是 Schwartz 函数。此时极限 $\lim_{f\to1}Q_f$ 不存在，但对任一局域泛函 $F$ 有
\begin{equation}
  \tag{1.19}
  \label{eq:1-19}
  \lim_{f\to1}\{F,Q_f\}=\left.\frac{d}{ds}\right|_{s=0}(G^s)^*F.
\end{equation}

\textbf{注（Remark）.} 所谓 $f\to1$，是指 $f$ 在任一紧集上一致收敛到 1，且 $f$ 的所有导数在整个空间上一致趋于零。

在量子理论中，故事相同，只是（局域）泛函被换为局域算符，Poisson 括号被换为对易子乘以 $i$。也就是说，对任一单参数对称性 $G^s$，存在取值于 $(d-1)$-形式的局域算符 $J(x)$，它守恒；并且电荷算符 $Q=\int_C J(x)$（在有限体积下）具有性质
\begin{equation}
  \tag{1.20}
  \label{eq:1-20}
  [F,Q]=-i\left.\frac{d}{ds}\right|_{s=0}(G^s)^*F.
\end{equation}

% [Tutor Note: (1.19) 与 (1.20) 之间相差因子 $-i$，这对应 Dirac 量子化规则 $\{\,,\}\to\frac1i[\,,\ ]$；但符号要小心：若按 (1.19) 定义 $\{F,Q\}\equiv D$（$D=\frac{d}{ds}|_{s=0}(G^s)^*F$），标准规则给出的是 $[F,Q]=iD$，而 (1.20) 写的是 $-iD$。二者之差等价于把生成元取为 $G^s=e^{-isQ}$ 而不是 $e^{isQ}$，也就是 $\eta$（从而 $J$）的整体定向约定。由于后文只用到 $[Q_f,\phi]$，(1.21) 中这个整体符号无关紧要；但读者若拿 (1.20) 与教科书中 $[Q,\phi]=-i\,\delta\phi$ 的写法对照，务必先统一 $Q$ 的符号约定，否则会误以为 (1.19)–(1.20) 自相矛盾。另注意 $\lim_{f\to1}Q_f$ 不存在而 $\lim_{f\to1}\{F,Q_f\}$ 存在，原因是对易子把积分限制在 $F$ 的支集附近，这正是「无穷体积下电荷算符不是算符、但其导数是」的标准图像。最后，作者说「$\{\,,\}$ 是 $S$ 上的 Poisson 括号」，而全篇解空间记作 $X$，此处 $S$ 疑为 $X$ 之误（或指前文引入的相空间），译文保留原符号 $S$。]

无穷体积的情形与经典理论中同样处理，即考虑截断算符 $Q_f$。

现在我们来讨论 Goldstone 定理。设由 $\mathcal{L}$ 定义的理论中的对称性 $G^s$ 不仅在经典上破缺，而且在量子力学意义上也破缺。在此情形，若我们有理论的一个解 $\mathcal{H}$（一个带有算符代数作用的 Hilbert 空间），则存在一个标量局域算符 $\phi$，其单点函数在对称变换下并非不变。

\textbf{注（Remark）.} 严格说来，我们只知道某个 $n$ 点函数并非不变；但在一切已知的对称性破缺情形中，也存在一个并非不变的单点函数。（利用算符乘积展开可以论证：若存在一个并非不变的 $n$ 点函数，则存在一个并非不变的单点函数。）

设 $Q_f$ 是对称性 $G^s$ 的截断电荷算符。我们有
\begin{equation}
  \tag{1.21}
  \label{eq:1-21}
  \langle\Omega|[Q_f,\phi(0)]|\Omega\rangle\neq0,
\end{equation}
其中 $f$ 充分接近 1。于是对某个 $x$ 有 $\langle\Omega|[J(x),\phi(0)]|\Omega\rangle\neq0$。

% ===== 原文 p.20 =====
考虑两点函数
\begin{equation}
  \tag{1.22}
  \label{eq:1-22}
  M_+(x):=\langle\Omega|J(x)\phi(0)|\Omega\rangle,\qquad
  M_-(x):=\langle\Omega|\phi(0)J(x)|\Omega\rangle
\end{equation}
（它们是 $V$ 上的 $(d-1)$ 形式）。由于对称性破缺，$M_+\neq M_-$；不过由类空间隔，当 $x$ 类空时二者重合。

设 $\mathcal{H}=\int_{p\in V_+}\mathcal{H}_p$ 是 $\mathcal{H}$ 关于平移群作用的谱分解（这里 $V_+$ 是整个光锥的正向部分）。利用内积对中间态（intermediate state）的分解
\[
\begin{aligned}
\langle J(x)\Omega,\phi(0)\Omega\rangle
&=\sum_n\langle J(x)\Omega,n\rangle\langle n,\phi(0)\Omega\rangle,
\end{aligned}
\]
我们得到
\begin{equation}
  \tag{1.23}
  \label{eq:1-23}
  M_\pm(x)=\int_{V_+}K_\pm(x,p)\,dp,
\end{equation}
其中 $K_+(x,p)$、$K_-(x,p)$ 是四动量为 $p$ 的中间态对 $M_+$、$M_-$ 的贡献（$K_\pm$ 是 $p$ 的向量值分布）。

由于 Lorentz 不变性，分布 $K_+$、$K_-$ 形如
\begin{equation}
  \tag{1.24}
  \label{eq:1-24}
  K_\pm=p\,e^{\pm ipx}\rho_\pm(-p^2),
\end{equation}
其中 $\rho_\pm(s)$ 是半直线 $s\ge0$ 上的分布。

% [Tutor Note: 这条论证里有两个容易混淆的输入，作者只用一句话带过。第一，「$M(x)$ 对类空的 $x$ 为零」用的是局域算符在类空间隔上的对易性（因果性）：$[J(x),\phi(0)]=0$（$x$ 类空），故 $M_+=M_-$。第二，「$W_m(x)=W_m(-x)$（$x$ 类空）」用的是 Klein–Gordon 传播子对 $x\to-x$ 的偶性。两者是不同的性质，缺一不可：前者给 $M_+-M_-$ 的类空消失，后者才使 (1.26) 中的 $\rho_+-\rho_-$ 能进一步化为 (1.27) 的形式。复算时若把 $x\to-x$ 的偶性误用到类时区（那里 $W_m(x)\neq W_m(-x)$，差一个 $2i\,\mathrm{Im}$ 项），就会错失 (1.28) 的推论。]

% [Tutor Note: (1.24) 把指标「藏」起来了：$J$ 是 $(d-1)$ 形式，故 $K_\pm$ 实际带有矢量指标，准确写法是 $K_\pm{}^\mu(x,p)=p^\mu e^{\pm ipx}\rho_\pm(-p^2)$。Lorentz 不变性只允许 $K^\mu$ 正比于唯一可用的四维向量 $p^\mu$（其余标量因子只能是 $p^2$ 的函数），而之所以写成 $\rho_\pm(-p^2)$ 而非 $\rho_\pm(p^2)$：本讲采用 $(-,+,+,+)$ 号差（由 (1.25) 下方的 $p^2=-m^2$ 与类空的 $x^2>0$ 可见），质量壳上 $-p^2=m^2\ge0$，这样宗量才落在半直线 $s\ge0$ 上，与「$\rho_\pm(s)$ 定义在 $s\ge0$」相符。若把 $p^\mu$ 写成 $p_\mu$，或漏掉 $e^{\pm ipx}$ 中与 $\rho_\pm$ 相应的 $\pm$，(1.25) 里 $W_m(\pm x)$ 的 $\pm$ 就会配错。]

由此得到
\begin{equation}
  \tag{1.25}
  \label{eq:1-25}
  M_\pm(x)=i*d_V\int_0^\infty\rho_\pm(m^2)W_m(\pm x)\,dm^2,
\end{equation}
其中 $W_m(x)=\int_{\mathcal{O}^+_m}e^{ipx}dp$ 是 [I-Perturbative QFT] 第 1 讲中定义的、质量为 $m$ 的 Klein–Gordon 传播子（$\mathcal{O}^+_m$ 是双曲面 $p^2=-m^2$ 的上叶）。

令 $M(x)=M_+(x)-M_-(x)$。由于当 $x$ 类空时 $W_m(x)=W_m(-x)$，对类空的 $x$，(1.25) 给出
\begin{equation}
  \tag{1.26}
  \label{eq:1-26}
  M(x)=i*d_V\int_0^\infty\bigl(\rho_+(m^2)-\rho_-(m^2)\bigr)W_m(x)\,dm^2.
\end{equation}

然而，正如我们已经提到的，由类空间隔，$M(x)$ 对类空的 $x$ 为零。这意味着 $\rho_+=\rho_-=\rho$，于是 (1.25) 给出
\begin{equation}
  \tag{1.27}
  \label{eq:1-27}
  M(x)=i*d_V\int_0^\infty\rho(m^2)\bigl(W_m(x)-W_m(-x)\bigr)dm^2.
\end{equation}

在满足 $x^2>0$ 的点 $x$ 处对 (1.27) 两端求导，并利用流的守恒与 Klein–Gordon 方程 $\nabla^2W_m=-m^2W_m$，我们得到
\begin{equation}
  \tag{1.28}
  \label{eq:1-28}
  \int_0^\infty m^2\rho(m^2)\bigl(W_m(x)-W_m(-x)\bigr)dm^2=0.
\end{equation}

取 Fourier 变换，我们得到
\begin{equation}
  \tag{1.29}
  \label{eq:1-29}
  p^2\rho(-p^2)=0.
\end{equation}

% [Tutor Note: 从 (1.27) 到 (1.28) 的「differentiating both sides」跳过了三步：先对 (1.27) 取外微分，用流守恒 $d_VJ=0$（即 $\partial_\mu M^\mu=0$）把微分从 $d_V$ 转移到 $W_m$ 上；再用 Klein–Gordon 方程把 $W_m$ 的二阶导数换成 $-m^2W_m$，于是 $m^2$ 作为本征值被「提」到积分号内（这就是被积函数多出因子 $m^2$ 的来源）；最后约去公共的 $i$ 与 $d_V$。这里还有一个符号陷阱：作者在 (1.25) 下方写 $p^2=-m^2$、在本页写类空的 $x^2>0$，说明号差是 $(-,+,+,+)$；在该号差下 $\partial_\mu\partial^\mu e^{ipx}=-p^2e^{ipx}=m^2e^{ipx}$，故与 (1.28) 相容的写法应是 $\nabla^2W_m=+m^2W_m$，即他把 $\nabla^2$ 定义为 $-\partial_\mu\partial^\mu$。由于 (1.28) 右端为零，这个整体符号不影响结论。而 (1.29) 说明 $\rho$ 的支撑集在 $p^2=0$：$\langle\Omega|[J(x),\phi(0)]|\Omega\rangle\neq0$ 只能由无质量中间态贡献，这个无质量粒子即 Goldstone 玻色子（此即 Goldstone 定理的谱表述）。]
